% !TeX root = ../main.tex
% 6. fejezet — Hirdetési piactér
% Hivatkozott források: Backend/HUNting.API/Controllers/AdsController.cs, Services/AdService.cs,
%   Services/AdQueryOptions.cs, Models/PagedResult.cs, Models/Ads/*, Data/HuntingContext.cs,
%   Frontend/HUNting/src/services/adsApi.ts, pages/Ads/AdsListPage.tsx, AdCreatePage.tsx,
%   AdEditPage.tsx, AdDetailPage.tsx, AdForm.tsx

\chapter{Hirdetési piactér}
\label{chap:hirdetesek}

A hirdetési piactér a HUNting legnagyobb forgalmú, publikusan is elérhető funkcióterülete:
vadászati eszközök, felszerelések és vendégvadász-meghívások adásvételét támogatja. A fejezet a
funkciót a már ismertetett sablon szerint, az adatmodelltől a felhasználói felületig mutatja be,
külön kitérve a szerveroldali lapozás-szűrés és a képfeltöltés megvalósítására.

\section{A funkció célja és felhasználói működése}
\label{sec:ads-cel}

A piactér böngészése nem igényel bejelentkezést: a hirdetéslista kulcsszó szerint kereshető, és
kategória, valamint minimum- és maximumár szerint szűrhető; a találatok laponként kilencesével,
létrehozási idő szerint csökkenő sorrendben jelennek meg. A listaelemről a hirdetés részletező
oldalára lehet továbblépni, amely a leírást, az árat, a helyszínt, a feltöltött képeket és a
kapcsolatfelvételi adatokat mutatja. Bejelentkezett felhasználó új hirdetést adhat fel: az űrlapon
cím, leírás, kategória, ár, helyszín és elérhetőségek adhatók meg, a képek pedig húzd-és-ejtsd
módszerrel, vágólapról beillesztve vagy fájlböngészővel csatolhatók. A hirdetés tulajdonosa -- és
az adminisztrátor -- a részletező oldalról szerkesztheti és törölheti a hirdetést, valamint
egyenként távolíthatja el a képeket; bejelentkezett felhasználó a listaoldalon a „Saját
hirdetéseim” kapcsolóval csak a maga hirdetéseire szűkítheti a találatokat.

% TODO: SC-02 képernyőkép (figures/sc-02-ads-list.png) — aktív kategória- és árszűrővel.
\begin{figure}[htbp]
	\centering
	\fbox{\parbox[c][5cm][c]{0.85\textwidth}{\centering SC-02: A hirdetéslista szűrőkkel és lapozóval (elkészítendő)}}
	\caption{A hirdetési piactér listaoldala keresővel, kategória- és árszűrőkkel}
	\label{fig:ads-list-screen}
\end{figure}

% TODO: SC-03 képernyőkép (figures/sc-03-ad-detail.png) — tulajdonosként, a kezelőgombokkal.
\begin{figure}[htbp]
	\centering
	\fbox{\parbox[c][5cm][c]{0.85\textwidth}{\centering SC-03: A hirdetés részletező oldala (elkészítendő)}}
	\caption{Hirdetés részletei képekkel, kapcsolati adatokkal és tulajdonosi műveletekkel}
	\label{fig:ad-detail-screen}
\end{figure}

\section{Adatmodell}
\label{sec:ads-adatmodell}

A funkció négy entitásra épül (\texttt{Models/Ads/}). A központi \texttt{Ad} entitás tárolja a
hirdetés címét, leírását, árát (\texttt{decimal?}), a helyszín szöveges megnevezését és opcionális
koordinátáit, a kapcsolatfelvételi e-mail-címet és telefonszámot, az állapotkódot, valamint a
létrehozás időbélyegét. A kategóriákat az \texttt{AdCategory} szótártábla írja le, amelynek neve
egyedi indexszel védett; a kezdeti adatfeltöltés tizenkét kategóriát hoz létre (például
„Fegyverek”, „Lőszerek”, „Vadkamera”). A hirdetés állapotát az \texttt{AdStatus} referenciatábla
kódolja -- itt nem numerikus, hanem karakterláncos elsődleges kulcsot használ a rendszer
(\texttt{ACTIVE}, \texttt{PENDING}, \texttt{SOLD}, \texttt{HIDDEN}) --, így az állapotkód
önmagában is beszédes. A képeket az \texttt{AdImage} entitás kapcsolja a hirdetéshez: a tárolt
útvonal mellett \texttt{SortOrder} rendezőszámot is tartalmaz, és a hirdetés törlésekor
kaszkádolva törlődik.

A tulajdonlás a modell legérdekesebb megszorítása. A hirdetés gazdája kétféle lehet --
felhasználó (\texttt{OwnerUserId}) vagy társaság (\texttt{OwnerSocietyId}) --, és a
\texttt{HuntingContext} egy adatbázisszintű ellenőrző megszorítással
(\texttt{CK\_Ads\_SingleOwner}) garantálja, hogy a két idegen kulcs közül \emph{pontosan az egyik}
legyen kitöltve. Ugyanezt az invariánst az üzleti réteg is ellenőrzi (lásd
a~\ref{sec:ads-backend}.~alfejezetet), így a szabály megsértése már az adatbázis elérése előtt
kiderül. Mindkét tulajdonosi idegen kulcs kaszkád törlésű: a felhasználó vagy a társaság
törlésekor a hirdetései is törlődnek; a kategória- és státuszhivatkozás ezzel szemben
\texttt{Restrict} viselkedésű, tehát használatban lévő szótárérték nem törölhető.

\section{Backend-megvalósítás}
\label{sec:ads-backend}

A funkciót az \texttt{AdsController} és a mögötte álló \texttt{AdService} szolgáltatás valósítja
meg; a kontroller kizárólag a HTTP-réteggel és a jogosultság-ellenőrzéssel foglalkozik, a
lekérdezés-összeállítás és az adatműveletek a szolgáltatásban élnek. A végpontokat
a~\ref{tab:ads-endpoints}.~táblázat foglalja össze.

\begin{table}[htbp]
	\centering
	\caption{Az \texttt{AdsController} végpontjai}
	\label{tab:ads-endpoints}
	\begin{tabular}{|l|l|p{5.9cm}|}
		\hline
		\textbf{Végpont} & \textbf{Hozzáférés} & \textbf{Feladat} \\
		\hline
		GET \texttt{/api/Ads} & anonim & Lapozott, szűrhető hirdetéslista \\
		\hline
		GET \texttt{/api/Ads/categories} & anonim & Kategórialista \\
		\hline
		GET \texttt{/api/Ads/statuses} & anonim & Állapotlista \\
		\hline
		GET \texttt{/api/Ads/\{id\}} & anonim & Egy hirdetés részletei \\
		\hline
		POST \texttt{/api/Ads} & bejelentkezett & Hirdetés létrehozása \\
		\hline
		PUT \texttt{/api/Ads/\{id\}} & tulajdonos v. admin & Hirdetés módosítása \\
		\hline
		DELETE \texttt{/api/Ads/\{id\}} & tulajdonos v. admin & Hirdetés törlése \\
		\hline
		POST \texttt{/api/Ads/\{id\}/images} & tulajdonos v. admin & Kép feltöltése (max.\ 10~MiB) \\
		\hline
		DELETE \texttt{/api/Ads/\{id\}/images/\{imageId\}} & tulajdonos v. admin & Kép törlése \\
		\hline
	\end{tabular}
\end{table}

\subsection{Lapozott és szűrt lekérdezés}

A listázó végpont a lekérdezési paramétereket az \texttt{AdQueryOptions} osztályba köti
(oldalszám, lapméret, keresőkifejezés, kategória, állapotkód, ár-intervallum, tulajdonos). A
szolgáltatás a lapozási értékeket normalizálja -- egynél kisebb oldalszámból egy lesz, az 1 és 100
közé nem eső lapméret helyett az alapértelmezett tizenkettő érvényesül --, majd a szűrőfeltételeket
csak akkor fűzi a lekérdezéshez, ha az adott paraméter ténylegesen meg van adva
(\ref{lst:ads-paged}.~kódrészlet). Így minden szűrés az adatbázisban, SQL-szinten fut le, a
válasz pedig a \texttt{PagedResult<T>} rekordban a találatok mellett a teljes találatszámot és a
lapozási adatokat is visszaadja -- a kliens ebből számítja az oldalszámot.

\begin{lstlisting}[language={[Sharp]C}, caption={Szerveroldali lapozás és feltételes szűrés az \texttt{AdService.GetPagedAsync} metódusban (kivonat)}, label={lst:ads-paged}]
var page = options.Page < 1 ? 1 : options.Page;
var take = options.Take is < 1 or > 100 ? 12 : options.Take;

IQueryable<Ad> query = _context.Ads
    .AsNoTracking()
    .Include(a => a.Category)
    .Include(a => a.Images)
    .Include(a => a.Status);

if (!string.IsNullOrWhiteSpace(options.Search))
{
    var term = options.Search.Trim();
    query = query.Where(a =>
        EF.Functions.Like(a.Title, $"%{term}%") ||
        (a.Description != null && EF.Functions.Like(a.Description, $"%{term}%")));
}

if (options.MinPrice.HasValue)
{
    query = query.Where(a => a.Price >= options.MinPrice.Value);
}

var total = await query.CountAsync();

var items = await query
    .OrderByDescending(a => a.CreatedAt)
    .Skip((page - 1) * take)
    .Take(take)
    .ToListAsync();

return new PagedResult<Ad>(items, total, page, take);
\end{lstlisting}

A keresés az \texttt{EF.Functions.Like} függvénnyel fordul SQL \texttt{LIKE} kifejezésre, és a
címre, valamint a leírásra egyaránt illeszkedik. A lekérdezés \texttt{AsNoTracking} módban fut,
mivel az eredmény csak olvasásra szolgál, a kapcsolódó kategória-, kép-, állapot- és
tulajdonos-adatokat pedig \texttt{Include} hívások töltik be egyetlen lekérdezésben.

\subsection{Írási műveletek és tulajdonlás}

Hirdetés létrehozásakor a kontroller a hívó szerepétől függően kezeli a tulajdonlást: nem
adminisztrátor esetén a tulajdonos kényszerítetten maga a hívó (a kérésben esetleg megadott
tulajdonosmezők felülíródnak), adminisztrátor viszont tetszőleges felhasználói vagy társasági
tulajdonost adhat meg -- ha egyiket sem adja meg, a rendszer őt magát jelöli tulajdonosnak. Az
\texttt{AdService} minden írási művelet előtt a „pontosan egy tulajdonos” invariánst is ellenőrzi,
és megsértése esetén \texttt{ArgumentException} kivételt dob, amelyet a kontroller 400-as
válaszra fordít. A módosítás és a törlés a~\ref{chap:felhasznalok}.~fejezetben bemutatott
erőforrás-tulajdonos mintát követi (\ref{lst:owner-check}.~kódrészlet): a műveletet csak a
tulajdonos vagy adminisztrátor végezheti el, és a tulajdonosmezők átírása adminisztrátori
kiváltság.

\subsection{Képfeltöltés}

A képfeltöltő végpont többrészes (\texttt{multipart/form-data}) kérést fogad, amelynek méretét a
\texttt{[RequestSizeLimit]} attribútum 10~MiB-ban korlátozza. A feldolgozás több lépcsős
védelmet alkalmaz: üres fájl esetén azonnali hibával tér vissza, a fájlnév kiterjesztését
kisbetűsítve fehérlistával veti össze (csak \texttt{.jpg}, \texttt{.jpeg}, \texttt{.png} és
\texttt{.webp} engedélyezett), a lemezre pedig nem az eredeti fájlnéven, hanem újonnan generált
GUID néven ír -- ezzel kizárva az útvonal-befecskendezést és a fájlnév-ütközést. A fájl a webes
gyökér \texttt{uploads/ads} könyvtárába kerül (a könyvtárat a kontroller szükség esetén
létrehozza), az adatbázisba csak a relatív URL íródik, a kiszolgálást pedig a statikus
fájl-middleware végzi. Az \texttt{AdService.AddImageAsync} a rendezőszámot a meglévő képek
maximuma alapján, folytatólagosan osztja ki. Kép törlésekor a végpont ellenőrzi, hogy a kép
valóban az útvonalban megadott hirdetéshez tartozik-e, majd a lemezen tárolt fájlt és az
adatbázisrekordot is eltávolítja.

\section{Frontend-megvalósítás}
\label{sec:ads-frontend}

\subsection{Listaoldal és szűrés}

Az \texttt{AdsListPage} a szűrőket (keresőmező, kategória-legördülő, ár-intervallum, „Saját
hirdetéseim” kapcsoló) komponens-lokális állapotban tartja, a találatokat pedig a React Query
kezeli. A lekérdezés kulcsa a teljes szűrőállapotot tartalmazza
(\texttt{["ads", categoryId, search, minPrice, maxPrice, showMyAds, page]}), így bármely szűrő
változása automatikusan új lekérdezést vált ki, a korábbi szűrőkombinációk eredménye pedig a
gyorsítótárból azonnal visszatölthető. Egy effektus gondoskodik róla, hogy szűrőváltáskor a lapozó
az első oldalra ugorjon. A „Saját hirdetéseim” kapcsoló a bejelentkezett felhasználó azonosítóját
adja át \texttt{ownerUserId} paraméterként, vagyis a szűrés ilyenkor is szerveroldalon történik. A
lapszámot a kliens a válasz \texttt{total} mezőjéből számítja, kilences lapmérettel.

\subsection{Hirdetésfeladás és az űrlap}

Az új hirdetés oldala (\texttt{AdCreatePage}) vendég látogatót a bejelentkezési oldalra irányít,
az eredeti célútvonalat útvonalállapotban megőrizve. Az űrlapot a létrehozó és a szerkesztő oldal
által közösen használt \texttt{AdForm} komponens adja: kontrollált beviteli mezőkkel dolgozik, a
kategórialistát a gyorsítótárazott \texttt{["ad-categories"]} lekérdezésből tölti, az új képeket
pedig memóriában gyűjti -- a fájlok háromféleképpen adhatók hozzá (fájlböngésző, húzd-és-ejtsd,
vágólap-beillesztés), és csak \texttt{image/} MIME-típusú fájlokat fogad el. Az előnézetekhez a
komponens objektum-URL-eket képez, amelyeket leszereléskor fel is szabadít. Szerkesztéskor a már
feltöltött képek nem azonnal törlődnek, hanem eltávolításra jelölhetők, és csak az űrlap
beküldésekor kerülnek ténylegesen törlésre.

A beküldés folyamata kétfázisú, mert a backend a hirdetést és a képeket külön végpontokon
kezeli: először a hirdetés jön létre, majd a kapott azonosítóval egyenként töltődnek fel a képek,
végül a kliens érvényteleníti a hirdetéslista gyorsítótárát és a részletező oldalra navigál
(\ref{lst:ads-create}.~kódrészlet).

\begin{lstlisting}[caption={A hirdetésfeladás kétfázisú beküldése az \texttt{AdCreatePage} komponensben (kivonat)}, label={lst:ads-create}]
const created = await mutation.mutateAsync(payload);

if (newFiles.length > 0) {
  setIsProcessingImages(true);
  for (const file of newFiles) {
    await adsApi.uploadImage(created.adId, file);
  }
}

await queryClient.invalidateQueries({ queryKey: ["ads"] });
await queryClient.invalidateQueries({ queryKey: ["ads", created.adId] });
navigate(`/ads/${created.adId}`, { replace: true });
\end{lstlisting}

\subsection{Részletező és szerkesztő oldal}

A részletező oldal (\texttt{AdDetailPage}) az útvonal-paraméterből kiolvasott azonosítóval, az
\texttt{["ads", adId]} kulcson kérdezi le a hirdetést. A kezelőgombok megjelenítéséről a
\texttt{canManage} feltétel dönt: a felhasználó adminisztrátor, vagy a hirdetés
\texttt{ownerUserId} mezője egyezik az azonosítójával -- ez ugyanaz a szabály, amelyet a backend
is kikényszerít, így a felület nem ígér olyan műveletet, amelyet a szerver elutasítana. A törlés
natív megerősítő kérdés után fut, sikere a lista-gyorsítótár érvénytelenítésével és a listaoldalra
navigálással zárul. Képtörléskor a kliens optimalizált frissítést alkalmaz: a részletező oldal
gyorsítótárazott adatából azonnal kiszűri a törölt képet, és csak ezután érvényteleníti a
listákat. A szerkesztő oldal (\texttt{AdEditPage}) ugyanazt az \texttt{AdForm} komponenst tölti
fel a meglévő hirdetés adataival; beküldéskor a módosítás mellett az újonnan csatolt képeket
feltölti, az eltávolításra jelölteket törli, majd szintén a részletező oldalra tér vissza.

% TODO: SC-04 képernyőkép (figures/sc-04-ad-form.png) — képelőnézettel és drag&drop zónával.
\begin{figure}[htbp]
	\centering
	\fbox{\parbox[c][5cm][c]{0.85\textwidth}{\centering SC-04: A hirdetésfeladó űrlap (elkészítendő)}}
	\caption{A hirdetésfeladó űrlap képelőnézettel és húzd-és-ejtsd feltöltési zónával}
	\label{fig:ad-form-screen}
\end{figure}

\section{Adatfolyam, biztonság és korlátok}
\label{sec:ads-adatfolyam}

A hirdetésfeladás teljes, képfeltöltéssel kombinált adatfolyamát
a~\ref{fig:ad-create-sequence}.~ábra szemlélteti. A folyamat jól példázza a REST-erőforrások
szétválasztásából adódó tervezési kompromisszumot: a hirdetés és a képek külön erőforrások, ezért
a kliens felelőssége a műveletsor levezénylése. Előnye, hogy a képfeltöltés hibája nem veszíti el
magát a hirdetést; hátránya, hogy a folyamat nem atomi -- ha a feltöltési ciklus félbeszakad, a
hirdetés a képek egy részével marad meg.

% TODO: F-01 szekvenciadiagram (figures/f-01-ad-create-sequence.pdf) — AdForm → POST /ads →
%       AdService → SaveChanges → képenként POST /ads/{id}/images → wwwroot/uploads.
\begin{figure}[htbp]
	\centering
	\fbox{\parbox[c][6.5cm][c]{0.9\textwidth}{\centering F-01: A hirdetésfeladás szekvenciadiagramja képfeltöltéssel (elkészítendő)}}
	\caption{A hirdetésfeladás és képfeltöltés szekvenciadiagramja}
	\label{fig:ad-create-sequence}
\end{figure}

Biztonsági szempontból a funkció többrétegű védelmet valósít meg: a módosító végpontok
tulajdonos- vagy adminisztrátor-ellenőrzéshez kötöttek; a tulajdonlás kettős invariánsát az
üzleti réteg és adatbázis-megszorítás együtt őrzi; a képfeltöltést kiterjesztés-fehérlista,
méretkorlát és generált fájlnév védi. A korlátok közé tartozik, hogy a kiterjesztés-ellenőrzés a
fájlnéven alapul, a fájl tényleges tartalmát (mágikus bájtok) a rendszer nem vizsgálja; a
feltöltött fájlok autentikáció nélkül, statikus tartalomként érhetők el, ami nyilvános piactérnél
elfogadható, de érzékeny tartalomnál nem volna az. Megjegyzendő továbbá, hogy a végpontok magát az
EF-entitást fogadják és adják vissza kérés- és válaszmodellként; ez kevesebb kódot igényel, mint a
dedikált DTO-réteg, cserébe a módosító kérések a szükségesnél több mezőt tartalmazhatnak -- a
kontroller ezért másolja át mezőnként, explicit listával a módosítható értékeket. Végül a
kliensoldali állapotlista jelenleg rögzített \texttt{ACTIVE} alapértelmezettel dolgozik, és bár az
állapotlekérdező végpont és kliensfüggvény létezik, a vizsgált oldalak nem használják -- az
állapot alapú szűrés így a felületen kiaknázatlan lehetőség maradt.
